\section{Ykrul --- ``Wolf Standards, Winter Camps''}
\label{chap:ykrul}

\subsection*{Quick Reference}
\textbf{Tagline:} \textit{The steppe remembers every oath. The wolf counts your breaths.}

\textbf{Genre / Mood:} Nomadic epic, honour‑bound geometry, wind‑whispered law.

\textbf{Starting Location:} A winter camp on the steppe edge, banners snapping in frost‑bright wind as the Khagan's fires burn in the east. A Khatun of the Ring offers a clay cup of salted tea. ``The Khagan's kin ride from the salt pan. The old sayings say the straight road is a lie. You came by the bent path — the wolf knows you. Speak your errand before the wind changes.''

\begin{quote}
  \textbf{Khatun of the Ring (Elite):}  
  ``The steppe remembers every oath sworn beneath its endless sky, and the wolf‑standards carry the honour of clans that have outlasted empires. To lead here is to keep the old ways sharp as winter steel.''
\end{quote}

\begin{quote}
  \textbf{Herd-Scout (Commoner):}  
  ``The horses know the weather before it comes, and the wolves know the thoughts in a man's heart. Best speak true and ride steady – the steppe doesn't forgive lies or clumsy hands on the reins. An old saying: the straight road is a lie. The wise follow the bent twig.''
\end{quote}

\begin{tcolorbox}[colback=black!3,colframe=blue!60!black,title={Theme \& Atmosphere}]
The Ykrul ride the endless Violet Steppe, where the horizon drinks the sky and law lives in saddlebags. Their sacred patterns are called \textit{Kon'reh} — not a written code but a living way of reading routes, obligations, and exits. Every crossing has its price, every ford its weight, every feud an angle measured in horse-lengths and heartbeats. Hosts gather by the thousand, then scatter at a word; one night you share their fires, the next they test your bonds. Their banners are wolves' teeth, their treaties braided cords, their memory longer than the rivers that cut the grass. The old folk still recite the Thirty‑Six Sayings of the Road, but the wisest know the thirty‑seventh is the one you learn alone, on a night when the stars blink out and the wolf does not howl.

\vspace{2mm}
\textbf{What you notice first:} The smell of horse sweat and woodsmoke. The way every rider counts hoofbeats under their breath. The silence when a stranger approaches the fire.

\textbf{The one rule that kills newcomers:} Never ride between two wolves at dusk. The space between them is a blind angle — a route with no exit, a weight with no payment.
\end{tcolorbox}

\subsection*{Regional Mood: The Geometry of the Steppe}

Here, the horizon is not a line — it is a question. The Ykrul measure distance in oaths, not miles. A camp is a weave of favours owed and honour satisfied. A feud is a slow equation: blood in, blood out, with interest compounded by winter. The Khagan rules not by force alone but by keeping the old patterns — the unwritten rules that govern every crossing, every treaty, every duel. Break the pattern, and the steppe itself may turn against you. Keep it, and even enemies will share their salt. The Ykrul build no monuments. They build relationships, each one a stone in an invisible cairn that stretches from the Dolmis to the Aberderrin.

\subsection*{Prominent NPCs of the Ykrul}
\label{sec:ykrul-npcs}

\begin{longtable}{|p{0.2\textwidth}|p{0.25\textwidth}|p{0.3\textwidth}|p{0.15\textwidth}|}
\hline
\textbf{Name} & \textbf{Role/Title} & \textbf{Motivation} & \textbf{Rumoured Quirk} \\
\hline
Khatun Bortai & Keeper of the Ring & Keep the clans from tearing each other apart & Her hostage strings are braided from the hair of oath‑breakers. \\
\hline
Noyan Temur & Banner‑lord of the Red Reeds & Win a winter pasture without bloodshed & He has never lost a negotiation. He has never needed to speak. \\
\hline
Sogata & Herd‑Scout (commoner exemplar) & Read the weather before the storm & He can read a trail by the way the grass bends — and a lie by the way a rider sits. \\
\hline
Elder Saruul & Bone‑Singer & Keep the names of the dead & She sings the names of the departed each dawn; the dead sing back. \\
\hline
The Dust Prince & Raider who rides the wind's lee & Amass enough salt to buy a clan & He offers escort braids for secrets, never seen from the front. \\
\hline
Storm‑Singer Urum & Sky‑Speaker & Sell weather on credit & His bills come due during the Wolf Moon. \\
\hline
\end{longtable}

\subsection*{Names of the Ykrul}
\label{sec:ykrul-names}

Inspired by the open steppe and the deep east: guttural stops, flowing vowels, and suffixes that mark lineage or deed.

\begin{longtable}{|p{0.25\textwidth}|p{0.25\textwidth}|p{0.2\textwidth}|p{0.2\textwidth}|}
\hline
\textbf{First Names (Male)} & \textbf{First Names (Female)} & \textbf{Surnames} & \textbf{Epithets} \\
\hline
Bortai & Temur & Sogata & \textit{the Grey} \\
\hline
Urum & Saruul & Khorch & \textit{Red‑String} \\
\hline
Khasar & Altan & Noyan & \textit{the Unbuckled} \\
\hline
Mongke & Checheg & Borjigin & \textit{Storm‑Born} \\
\hline
Edil & Borte & Uriankhai & \textit{Ash‑Prow} \\
\hline
Toghrul & Orghana & Qasar & \textit{the Silent} \\
\hline
\end{longtable}

\paragraph*{(Steppe Road/Winter Ring/Ford)} Wolf Road milepost of stacked stones with wind whining through eye‑holes; winter camp ring with felt tents in horseshoe facing wind.

\section*{Spades --- Places}
\label{sec:ykrul-places}
\begin{enumerate}
\setcounter{enumi}{1}
\item \textbf{Wolf Milepost} --- Stacked stones; wind whines through eye‑holes. `[TRAVEL]`
  \hfill \textit{Each stone is a treaty. Remove one, and the wind names you oath‑breaker.}
\item \textbf{Remount Station} --- Low corral with stamped snow and tether‑posts. `[TRAVEL]`
  \hfill \textit{The tether‑posts are carved with the names of dead horses. The living ones remember.}
\item \textbf{Birch Windbreak} --- Above black‑earth tracks; old offerings in the bark. `[RITUAL]`
  \hfill \textit{The offerings are strips of wool. Each strip is a prayer for safe crossing.}
\item \textbf{Salt Pan} --- Crusted white; hoof‑prints like stars across the void. `[WEIGHT]`
  \hfill \textit{The pan has a name: ``The Khagan's Mirror.'' Look into it, and you see what you owe.}
\item \textbf{Reed Ford} --- River braids; poles mark yesterday's safe line. `[TRAVEL]`
  \hfill \textit{The poles shift in the night. The river's path changes with the moon.}
\item \textbf{Trading Palisade} --- Way‑post with wolf‑skull pennons snapping. `[TRADE]`
  \hfill \textit{The skulls face outward. They watch for liars and bad weights.}
\item \textbf{Winter Camp} --- Felt tents in a horseshoe facing the wind. `[SOCIAL]`
  \hfill \textit{The open end is for guests. The closed end is for the ancestors.}
\item \textbf{Kurgan Field} --- Low barrows; horse‑bones bead the grass like prayer beads. `[DEATH]`
  \hfill \textit{The bones were placed by a bone‑singer. Each bead is a name.}
\item \textbf{Watch Kopje} --- Lonely tor with signal fire‑scar; the eye of the steppe. `[WEATHER]`
  \hfill \textit{The fire‑scar is from a signal that saved a host. The smoke still smells of hope.}
\item[J] \textbf{Pontoon Crossing} --- Laced hides creak; the current speaks fast and hungry. `[TRAVEL]`
  \hfill \textit{The hides are from horses that died crossing. Their ghosts still pull.}
\item[Q] \textbf{Council Hollow} --- Standards planted; the ground tamped like a war‑drum. `[SOCIAL]`
  \hfill \textit{The standards have ears. Speak a treason here, and every clan will know by dawn.}
\item[K] \textbf{Khagan's Station} --- Broad ger on a timber platform; guards in silence. `[AUTHORITY]`
  \hfill \textit{The guards never blink. They are paid in salt and secrets.}
\item[A] \textbf{Sky Steppe} --- Open, wind like a voice; tracks keep pace beside you. `[MYSTERY]`
  \hfill \textit{The tracks are not yours. They belong to the ancestors who ride with you.}
\end{enumerate}

\paragraph*{(Host/Envoy/Rider)} Herd‑scout with three ponies and six opinions; camp‑mother who chooses where fires live.

\section*{Hearts --- People \& Factions}
\label{sec:ykrul-people}
\begin{enumerate}
\setcounter{enumi}{1}
\item \textbf{Herd-Scout} --- Three ponies and six opinions; the horizon as a map. `[TRAVEL]`
  \hfill \textit{He can read a trail by the way the grass bends. He can read a lie by the way a rider sits.}
\item \textbf{Camp-Mother} --- Chooses where fires live; hospitality is law. `[SOCIAL]`
  \hfill \textit{She has never drawn a blade. She has never needed to.}
\item \textbf{Banner Youth} --- Minding the wolves' tooth standard; pride as duty. `[COMBAT]`
  \hfill \textit{The standard is older than his grandmother's grandmother. He will die before it falls.}
\item \textbf{Salt-Broker} --- Weighs promises like grain; trade as precision. `[TRADE]`
  \hfill \textit{Her scales are honest---except when they aren't. The difference is a prayer.}
\item \textbf{Remount Keeper} --- String‑master with a book of brands; horses as currency. `[TRAVEL]`
  \hfill \textit{His book is written in hoof‑prints. He can read a horse's soul.}
\item \textbf{Bone-Singer} --- Knows the kurgans' proper names; death as memory. `[DEATH]`
  \hfill \textit{She sings the names of the dead each dawn. The dead sing back.}
\item \textbf{Road-Judge} --- Tümen scribe with seals and a patient ear. `[LAW]`
  \hfill \textit{His seals are from three empires. He uses whichever one is still valid.}
\item \textbf{Noyan Envoy} --- Silver paiza, iron smile; diplomacy as a blade. `[SOCIAL]`
  \hfill \textit{His paiza was a gift from a Khagan who is now dead. He does not mention that.}
\item \textbf{Host Captain} --- The tent‑wall is a map; strategy in every fold. `[COMBAT]`
  \hfill \textit{He can see a battle in the way smoke rises. He is usually right.}
\item[J] \textbf{Falcon Courier} --- Leather gauntlets, answers sooner than asked. `[TRAVEL]`
  \hfill \textit{The falcon is a sleeper agent. It belongs to a different clan. The courier doesn't know.}
\item[Q] \textbf{Khatun of Ring} --- Keeper of camp‑law and hostage strings. `[LAW]`
  \hfill \textit{Her strings are braided from the hair of hostages. She knows each strand's owner.}
\item[K] \textbf{Khagan's Kin} --- A nephew or niece sent to bind a foedus with a look. `[SOCIAL]`
  \hfill \textit{They have never lost a negotiation. They have never needed to speak.}
\item[A] \textbf{Sky-Speaker} --- Shaman of storms; shortcuts owe them a price. `[MAGIC]`
  \hfill \textit{They can call lightning---or call it off. The price is always a story.}
\end{enumerate}

\paragraph*{(Law/Weather/Feud)} White squall with dry snow, lost horizons; rasputitsa turning the road to porridge; a feud spark from a name spoken wrong.

\section*{Clubs --- Complications/Threats}
\label{sec:ykrul-complications}
\begin{enumerate}
\setcounter{enumi}{1}
\item \textbf{White Squall} --- Dry snow, lost horizons; navigation by instinct alone. `[WEATHER]`
  \hfill \textit{The snow tastes of iron. Someone is bleeding nearby.}
\item \textbf{Rasputitsa} --- The road becomes porridge; wheels are lies on the earth. `[WEATHER]`
  \hfill \textit{The mud has a smell---old blood and older lies.}
\item \textbf{Remount Sickness} --- A cough in the string; swaps get awkward, trust thinner. `[DEATH]`
  \hfill \textit{The cough is a curse. The horses know who cast it.}
\item \textbf{Salt Shortage} --- Meat spoils; tempers thin like winter air. `[WEIGHT]`
  \hfill \textit{The salt was cut with grave dust. Someone wanted the winter to be hard.}
\item \textbf{Hostage Call} --- Protocol invoked; someone must sit in the tent. `[SOCIAL]`
  \hfill \textit{The tent has no floor. The hostage can see the earth beneath---and the bones.}
\item \textbf{Feud Spark} --- A name spoken wrong reopens an old cut; blood remembers. `[WEIGHT]`
  \hfill \textit{The name was spelled correctly. The insult was in the pause before it.}
\item \textbf{Grassfire} --- Runs with the wind faster than horses; escape as a prayer. `[DISASTER]`
  \hfill \textit{The fire has a face. It is the face of a Khagan who burned a village.}
\item \textbf{Foedus Recall} --- Treaty terms re‑read; your papers disagree. `[LAW]`
  \hfill \textit{The treaty was written on both sides. The other side used a different ink.}
\item \textbf{River Break} --- Ice goes; ferries don't; crossing becomes a gamble. `[TRAVEL]`
  \hfill \textit{The ice is still thick in the middle. The edges are lying.}
\item[J] \textbf{Raid Shadow} --- Another banner follows at a patient distance. `[COMBAT]`
  \hfill \textit{They are not attacking. They are waiting for you to make a mistake.}
\item[Q] \textbf{Kurultai Call} --- Camps converge; travel becomes politics. `[SOCIAL]`
  \hfill \textit{The kurultai has a list of grievances. Your name is near the top.}
\item[K] \textbf{Wolf Muster} --- A levy raised; every road a checkpoint, every rider a judge. `[AUTHORITY]`
  \hfill \textit{The muster horn is a wolf's skull. It sounds like a scream.}
\item[A] \textbf{Sky Omen} --- Three suns, or none; routes rewrite themselves. `[OMEN]`
  \hfill \textit{The omen is a wedge of cranes flying north in winter. The sky is wrong.}
\end{enumerate}

\paragraph*{(Pass/Remount/Truce)} Camp token for one night's lawful fire and fodder; salt allotment from a named store without quarrel.

\section*{Diamonds --- Rewards/Leverage}
\label{sec:ykrul-rewards}
\begin{enumerate}
\setcounter{enumi}{1}
\item \textbf{Camp Token} --- One night's lawful fire and fodder; hospitality as a right. `[SOCIAL]`
  \hfill \textit{The token is a horsehair braid. Burn it to claim your fire.}
\item \textbf{Salt Allotment} --- Draw from a named store without quarrel. `[TRADE]`
  \hfill \textit{The allotment is a salt‑cake. Break it in half to share.}
\item \textbf{Ford-Right} --- Cross a named braid at lawful depth. `[TRAVEL]`
  \hfill \textit{The right is a pole. Plant it in the river, and the ford will hold---once.}
\item \textbf{Remount Chit} --- Swap two tired mounts for fresh. `[TRAVEL]`
  \hfill \textit{The chit is a horse's tooth. The remount keeper will not smile when you present it.}
\item \textbf{Escort Braid} --- Two riders see you through a bad stretch. `[TRAVEL]`
  \hfill \textit{The braid is woven from the riders' own hair. They will die before they cut it.}
\item \textbf{Safe-Hostage} --- Place or take token‑kin for a truce. `[SOCIAL]`
  \hfill \textit{The hostage is a child. The child is patient.}
\item \textbf{Paiza Tablet} --- Priority on the Wolf Road and at posts. `[AUTHORITY]`
  \hfill \textit{The tablet is bronze. It is warm when your reputation is good.}
\item \textbf{Foedus Seal} --- A treaty clause that trumps petty order. `[LAW]`
  \hfill \textit{The seal is a wolf's head. Press it into wax, and the steppe listens.}
\item \textbf{Market-Green} --- Trade under a banner; weapons sheathed. `[TRADE]`
  \hfill \textit{The banner is a green field. No blood may be spilled beneath it.}
\item[J] \textbf{Standard Guard} --- Travel under wolf‑tooth; raiders defer. `[COMBAT]`
  \hfill \textit{The standard is a wolf's tooth on a pole. Raiders will bow---or die.}
\item[Q] \textbf{Ring Audience} --- The Khatun hears you alone; precedent sticks. `[SOCIAL]`
  \hfill \textit{The audience is held at midnight. The fire is the only witness.}
\item[K] \textbf{Khagan's Writ} --- Doors open, mouths close across the steppe. `[AUTHORITY]`
  \hfill \textit{The writ is a sheet of hammered silver. It bends, but it does not break.}
\item[A] \textbf{Sky's Exception} --- Once, the weather ignores you (storm, whiteout, dust). `[WEATHER]`
  \hfill \textit{The exception is a single breath. Use it well.}
\end{enumerate}

\section*{Quick use notes}
\label{sec:ykrul-quick-use}
\begin{itemize}
\item Draw until you have all four suits: Spade = place, Heart = actor, Club = pressure, Diamond = leverage. Highest rank sets the main Clock (2–5 $\rightarrow$ 4, 6–10 $\rightarrow$ 6, J/Q/K $\rightarrow$ 8, A $\rightarrow$ 10).
\item Diamonds are codified outcomes (passes/remounts/truces) that change position rather than call for a roll.
\item If any A appears, echo \textbf{steppe and sky} motifs—wind that speaks, horses that remember, treaties that bind across generations.
\end{itemize}

\subsection*{Additional Features (Cultural Mechanics)}
\label{sec:ykrul-mechanics}

\begin{itemize}
\item \textbf{Hostage Strings.} Every oath is secured with a \textbf{hostage string}---a token kin, retainer, or sworn ally left in the other camp. If the deal is broken, the hostage's fate becomes the story's next beat. (Mechanically: create a 4‑segment \emph{Hostage} clock; on breach, the clock becomes a countdown to vengeance.)
\item \textbf{Steppe Law.} Travelers are bound by \textbf{steppe law}: hospitality lasts three nights, insults last three generations. The Keeper may demand 1 SB when either law is broken.
\item \textbf{Sky Omens.} When the sky shows strange signs---mock‑suns, green fire, thunder without storm---the steppe itself speaks. Once per scene, players may treat an omen as a \emph{Diamond}, shifting Position without a roll. The GM may bank 1 SB for future weather complications.
\item \textbf{Bowl and Board.} Ykrul elders settle disputes not by shouting but by drawing two circles in the dirt: the \textbf{Bowl} for fairness, then the \textbf{Board} for routes and concessions. A successful \emph{Wits + Lore} (or a relevant skill) roll creates a \emph{Seasonal Concession String} at a ford, harbor, or pass. On a miss, the \textbf{Market Moot} [4] clock advances as factions argue.
\item \textbf{Local Patrons (Syncretic Names):}
  \begin{itemize}
    \item \textbf{The Hoof} — The Traveler. \textit{The path that never ends, the remount that never tires.}
    \item \textbf{The Wolf's Tooth} — The Witness. \textit{Every oath is counted. Every broken promise is a scar.}
    \item \textbf{The Salt Throat} — Maelstraeus. \textit{The weight that salts your tongue; the interest that blinds your horses.}
    \item \textbf{The Sky Father} — Kuva. \textit{The wind that carries names; the storm that erases them.}
  \end{itemize}
\end{itemize}

\subsection*{The Thirty‑Seven Sayings of the Road (Excerpts)}
\label{sec:ykrul-sayings}

The Ykrul teach strategy not in manuals but in proverbs—each a line of play. A few of the most quoted:

\begin{quote}
\textit{Let the river freeze before you cross.} --- Do not rush. Let the opponent make the first move; your opening will be built on their impatience.

\textit{Borrow a corridor to trap a king.} --- When the enemy carves a path toward you, let it deepen. Then seal the mouth behind them.

\textit{Loot a burning house.} --- Strike only when the enemy is committed elsewhere.

\textit{Hide your dagger behind a smile.} --- Threaten one flank; strike the other.

\textit{The reed bends, the oak breaks.} --- Yield to overwhelming force, but do not break; preserve your core.

\textit{A distant banner casts no shadow.} --- A threat far away is no threat at all; control the centre.

\textit{The guest who outstays his welcome steals the host's peace.} --- Allow no force, however small, to remain within your lines uncontrolled.

\textit{Count exits, not victims.} --- Before you move, count the ways your foe can flee. If the number is greater than zero, you have not yet won.
\end{quote}

\subsection*{Lore Echoes (Cross‑Expansion Hooks)}
\begin{itemize}
  \item \textbf{The Bitter Root:} Daughters of the Cord run a hospice near the salt pans. They ask no payment---but they cut threads. A renegade Sister hides among the herders, trying to untie her cord from a Ykrul blood‑weight.
  \item \textbf{The Threshold Compendium:} The kurgans are Class II Anchor sites---the spirits of old Khagans bound by unfinished wergild. Saikou Ira would carry a bag of salt in each hand here.
  \item \textbf{Malachai, the Chained Angel:} A silver coin that never spends circulates among the salt‑brokers. Take it, and your store never runs dry---but the coin will reappear in the salt pan at every new moon.
  \item \textbf{The Grey Wanderer's Grimoire:} Ykrul psions are called ``Weather‑Readers.'' They taste the future in the wind. Road‑judges employ two---and have never trusted a verdict they didn't hear from the sky.
  \item \textbf{The Ninth:} In any Ykrul camp, the ninth cup poured at a feast is for the ancestors. The wise pour eight and leave the ninth as a pause.
\end{itemize}

\subsection*{Ykrul Superstitions (burn for +1 die once)}
\begin{itemize}
  \item \textbf{Salt to the Wind:} Toss a pinch into the air before a long ride. Ignore the first \emph{White Squall} penalty.
  \item \textbf{Knot the String:} Tie a single knot in a hostage cord. Cancel \emph{Feud Spark} or take –1 die on your next oath (your choice).
  \item \textbf{Name to the Kurgan:} Whisper a dead enemy's name into a horse's ear. Gain +1 Effect when raiding their clan this scene, but mark \emph{Obligation} to the Bone‑Singer.
\end{itemize}

\subsection*{Thematic SB Spend Table}
\label{sec:ykrul-sb}

\paragraph{Minor Complications (1 SB)}
\begin{itemize}
\item \textbf{Exposure:} Unwanted attention from \textbf{road‑judges or border riders}.
\item \textbf{Noise:} Sounds alert a \textbf{nearby camp or raiding party}.
\item \textbf{Trace:} Hoofprints or broken grass mark your route for \textbf{feud‑kin}.
\item \textbf{Delay:} A horse throws a shoe, a rope breaks — lose time.
\item \textbf{Supply Strain:} Mark +1 segment on \textbf{fodder or water} clock.
\end{itemize}

\paragraph{Moderate Setbacks (2 SB)}
\begin{itemize}
\item \textbf{Alarm Raised:} A \textbf{watch kopje} signals; riders are coming.
\item \textbf{Position Lost:} \textbf{Rasputitsa} or \textbf{white squall} hides the trail.
\item \textbf{Foe Appears:} A \textbf{raid shadow} or \textbf{hostage taker} arrives.
\item \textbf{Gear Trouble:} A saddle snaps, a bowstring frays.
\item \textbf{Lock/Barrier:} A river breaks, a ford becomes impassable.
\end{itemize}

\paragraph{Serious Trouble (3 SB)}
\begin{itemize}
\item \textbf{Reinforcements:} A \textbf{wolf muster} or \textbf{kurultai} war party arrives.
\item \textbf{Key Gear Breaks:} A horse goes lame, a paiza cracks.
\item \textbf{Major Twist:} A \textbf{grassfire} cuts off retreat, or a \textbf{feud spark} escalates.
\item \textbf{Rail Tick:} Advance a campaign clock (e.g. \textbf{Hostage Crisis}).
\item \textbf{Condition Applied:} Mark \textbf{Exhausted} or \textbf{Marked by Feud}.
\end{itemize}

\paragraph{Major Turns (4+ SB)}
\begin{itemize}
\item \textbf{Trap Springs:} A \textbf{hidden kurgan} collapses, a \textbf{pontoon} sinks.
\item \textbf{Authority Arrival:} \textbf{Khatun of the Ring}, \textbf{Khagan's Kin}, or \textbf{Sky-Speaker} intervenes.
\item \textbf{Scene Shift:} A \textbf{sky omen} rewrites the terrain; paths appear or vanish.
\item \textbf{Patron Omen:} A wolf watches, a horse speaks, the wind carries a name.
\item \textbf{Narrative Pivot:} The feud is not what it seemed; the hostage was a willing sacrifice.
\end{itemize}

\paragraph{Region-Specific SB Options}
\begin{itemize}
\item \textbf{Steppe Law:} Hospitality fails, insults echo, treaties shift mid‑parley.
\item \textbf{Horse Culture:} Mounts refuse to move, remount posts are empty, riding rolls suffer –1 die.
\item \textbf{Sky Omens:} Weather changes instantly, stars point the wrong way, a raven shadows the party.
\end{itemize}

\subsection*{Extensions (Plug‑in)}

\subsubsection*{Pursuit Ladder (Near / Mid / Far)}
\begin{itemize}
\item \textbf{Advance/Withdraw (Wits + Ride):} On a hit, shift one band. On a strong hit, also impose \emph{Dust Cloud}.
\item \textbf{Cut the Herd (Presence + Sway):} On a hit, freeze range for one exchange; if a \emph{remount chit} is in your possession, +1 die.
\item \textbf{Weather the Squall (Resolve + Survival):} On a hit in \emph{White Squall}, you pick who keeps their bearings; on a miss, the GM does.
\end{itemize}

\subsubsection*{Terrain Stock (d20)}
\label{sec:ykrul-terrain-stock}
Roll or choose a striking detail to dress a hex of the Violet Steppe.
\begin{enumerate}
  \item Wolf milepost with fresh red thread; someone retied the law.
  \item Reed ford poles lying on bank; current whispers a different depth.
  \item Fallen standard tooth used as a chopping block; grooves still warm.
  \item Circle of pony skulls facing east; one blinks when the wind shifts.
  \item Snow trench with cudgel steps; no prints leave it.
  \item Copper salt-scale hung from a birch; pans balance only under oath.
  \item Hoof-bells threaded on hair; ringing calls riders – or storms.
  \item Abandoned paiza shard wedged in a milestone; edges bite skin.
  \item Kurgan capstone chalked with a name misspelled; feud waiting.
  \item Flint knapper's tarp half-buried; a perfect arrowhead points north.
  \item Blue-fletched signal arrow stuck in frost; message in the vane.
  \item Camp-sweep marks: three rings – last night was a council.
  \item Wheel rut frozen like glass; a face looks up from inside.
  \item Salt-brined jerky hanging from a tripod; untouched by birds.
  \item River ice singing in thirds; ferryman's pole wedged under.
  \item Pony hobble-rope knotted in hostage-string pattern.
  \item Cold forge dug into loess; tools wrapped in wolfskin.
  \item Child's braid-bead dropped in snow; still warm to the tongue.
  \item Thorn-tether scratched with clan brands; two over-scratched.
  \item Bone flute lodged in a cairn; playing it summons a road-judge.
\end{enumerate}

\subsubsection*{Lair Seeds (One‑Page Rides)}
\label{sec:ykrul-lair-seeds}
\begin{enumerate}
  \item \textbf{Ger of Echoes} — Empty felt walls repeat anything uttered in the Khagan's voice once per night.
  \item \textbf{Salt-Pit Shrine} — Rope winch to brine caverns; offerings are tears and horsehair.
  \item \textbf{Remount Vault} — Subterranean stables under reed mats; mounts won't leave without a tale.
  \item \textbf{Kite-Maker's Bluff} — Bone kites pinned on racks; their shadows scout for raiders.
  \item \textbf{Storm-Rider's Stone} — Pillar that hums to hooves; racing round it sets weather wagers.
  \item \textbf{Frozen Court} — Ice dais in a dry riverbed; law held here counts triple until thaw.
  \item \textbf{Banner Grave} — Flags planted like trees; each knows the battle it died in.
  \item \textbf{String-Mother's Yard} — Web of hostage cords; cutting one opens a secret but starts a feud.
  \item \textbf{Fire-Under-Snow} — Warm vent where foals winter; spirits ask for salt-song.
  \item \textbf{Quiet Palisade} — Trading post where voices fail; deals done in gesture and cut meat.
  \item \textbf{Ashen Trident} — Three lightning-split birches; tie a braid to choose a road's luck.
  \item \textbf{Grey Ger} — Moves at night; inside, a camp-mother tests your hospitality the other way around.
\end{enumerate}

\subsubsection*{Steppe Bestiary (Palette)}
\label{sec:ykrul-bestiary}
\begin{itemize}
  \item \textbf{Steppe Manticores} — Lionine torsos, porcupine tails of bone quills. \emph{Signs:} quills in fence-posts; ponies refuse to drink. \emph{Hook:} play a raid-song backward; they circle to listen.
  \item \textbf{Dust Hulks} — Wind-packed silhouettes that rise from dunes. \emph{Signs:} dunes ``breathe''; grit sifts upward. \emph{Hook:} spill water on their shadow; they settle for a watch.
  \item \textbf{Bone Kites} — Carrion birds with rune-etched pinions. \emph{Signs:} spiral flights over kurgans. \emph{Hook:} offer a name-stake (twig with hair) to barter silence.
  \item \textbf{Frost Wargs} — Blue-pelt wolves that steam cold. \emph{Signs:} rime on tracks; breath that crackles. \emph{Hook:} warm milk poured on ice buys one unchallenged crossing.
  \item \textbf{Reed Stalkers} — Long-necked marsh hunters, gait like stilts. \emph{Signs:} pole‑pocks where no ferryman walks. \emph{Hook:} carry reeds in your mouth; they read it as truce.
  \item \textbf{Salt Screamers} — Pale eels in dry pans; voices like flutes. \emph{Signs:} fissures crusted with salt frost. \emph{Hook:} scatter ash; they coil to sing and forget to bite.
\end{itemize}

\begin{tcolorbox}[colback=black!3,colframe=black!40!white,title={Relic \& Find (d10)}]
\label{sec:ykrul-relic-find}
When you uncover a kurgan cache or a lost hoard, roll or pick:
\begin{enumerate}
  \item Wolf-tooth standard tip; planted, it demands silence for parley.
  \item Bronze paiza half; merchants honour it as full once – then talk.
  \item Salt-needle that points to the nearest brine store.
  \item Frost-bitten stirrup; mount that wears it ignores ice once.
  \item Braid-bead of hostage glass; melts when a string is cut nearby.
  \item River-depth rod carved with safe counts for three fords.
  \item Iron hoof-bell that rings for liars at council.
  \item Bone-lot cup; roll it to name which feud clause applies.
  \item Wind-wrung scarf; pass it round to share breath in a white squall.
  \item Wolf-road tablet scrap; reads as a \emph{Foedus Seal} to one clerk.
\end{enumerate}
\end{tcolorbox}

\subsection*{Complication Twists (d8)}
\label{sec:ykrul-twists}
When a travel complication arises, roll or choose a twist that deepens the story.
\begin{enumerate}
  \item \emph{White Squall} reveals a remount vault's smoke hole; entering angers reed stalkers.
  \item \emph{Rasputitsa} exposes kurgan stones; riding around them starts a name-claim.
  \item \emph{Feud Spark} was staged by a salt-broker; paying in \emph{Salt Allotment} unmasks them.
  \item \emph{Hostage Call} is for a \emph{Safe-Hostage} you already carry – by another name.
  \item \emph{Grassfire} runs uphill toward a council hollow; smoke votes before riders do.
  \item \emph{Foedus Recall} cites a river you crossed yesterday – now it flows the other way.
  \item \emph{Raid Shadow} is your own tracks, mirrored by sky omens; choose which you admit.
  \item \emph{Kurultai Call} stacks councils; words spoken in one bind you in the other.
\end{enumerate}

\subsection*{Boss Seeds (Three‑Phase Foes)}
\label{sec:ykrul-bosses}
\begin{itemize}
\item \textbf{The String-Mother (Nine Hostages)} — \emph{Phases}: Knot $\rightarrow$ Tighten $\rightarrow$ Sever. \emph{Cracks}: return a hostage string, win a Bowl and Board arbitration, cut the right knot.
\item \textbf{The Dust Prince (Rider of the Red Wind)} — \emph{Phases}: Raid Shadow $\rightarrow$ Salt Tithe $\rightarrow$ Kurultai Coup. \emph{Cracks}: spill a salt allotment in his path, unmask his face before the council, outride him through a white squall.
\end{itemize}

\begin{tcolorbox}[colback=black!3,colframe=black!40!white,title={Patronage \& Power}]
In Ykrul society, power flows through the careful management of hospitality, the binding of oaths with hostages, and the ability to navigate the complex web of tribal relationships. The Khagan maintains authority through the loyalty of noyans and the respect of khatuns, while local leaders wield influence through their control of resources like salt, remounts, and safe passage. The true power lies with those who understand the ancient customs and can maintain the delicate balance between hospitality and honour.

\textbf{For the GM:}  
Patronage in Ykrul society revolves around hospitality, hostage exchanges, and the ability to command respect from hosts and courts. Rewards often take the form of tokens, treaties, and official positions that can be leveraged into greater influence. To emphasize the old ways:
\begin{itemize}
\item Tie rewards to visible symbols (tokens, braids, standards) that can be challenged, stolen, or voided.
\item Use the \emph{Bowl and Board} procedure for disputes---first fairness (Bowl), then routes and concessions (Board), then a seasonal right.
\item Let rival hosts issue conflicting protections, forcing players to choose whose favour matters more.
\end{itemize}
In Ykrul society, your word is your bond, and your bond determines whether you ride as guest or ghost.
\end{tcolorbox}

% Index entries
\index{Ykrul|see{Wolf Standards, Winter Camps}}
\index{Theme generator}
\index{Wolf Standards, Winter Camps generator}
\index{Places!Ykrul}
\index{People!Ykrul}
\index{Complications!Ykrul}
\index{Rewards!Ykrul}
\index{Kon'reh}
\index{Thirty-Six Sayings}
\index{Hostage string}
\index{Steppe law}
\index{Sky omen}
\index{Bowl and Board}
\index{Local Patrons!Ykrul}

\subsection*{Ykrul of the Eastern Steppe — The Iron Grass Clans}
\label{sec:ykrul-east}

\emph{The wind is colder here, and so are the hearts. We do not forget an insult because we cannot afford to. The grass is thinner, the wolves are bolder, and the sky speaks only in threats. The western Ykrul trade in stories. We trade in teeth.}

\begin{quote}
  \textbf{Eastern Khatun (Elite):}  
  ``The westerners count their weights in silk and verses. We count ours in frost and iron. A broken oath here is not a story for the skald — it is a hole in the winter larder. We do not forgive because we cannot. The Hollow rides with us every night. We have no room for mercy.''
\end{quote}

The Eastern Ykrul share blood with their Violet Steppe cousins, but generations on the harsher eastern ranges — where the Kuvani horse clans also roam — have carved them into a different shape. They are more isolationist, more pragmatic, and their honour code has less room for poetry. They live among the Kuvani but do not intermarry; they trade when the snows allow, raid when the grass is thin, and speak a dialect that grinds like stones in a riverbed. Outsiders call them the \textit{Iron Grass Clans}. They call themselves \textit{Those Who Remember Winter}.

\subsubsection*{Eastern Variations (Use with Ykrul Core)}
When generating an Eastern Ykrul camp, clan, or encounter, roll on the tables below instead of (or in addition to) the standard Ykrul generators. Use dice rather than cards.

\paragraph*{Eastern Clan Trait (d6)}
Roll to determine a clan’s defining characteristic.
\begin{enumerate}
  \item \textbf{Frost-Bitten:} The clan survives on the edge of the permafrost. +1 die to Survival in cold weather, but –1 die to all social rolls with lowlanders (they stink of frozen horse).
  \item \textbf{Kuvani-Kin:} This clan has traded with the Kuvani for generations. They speak the Kuvani tongue and may call on one Kuvani remount station per journey. Rivals whisper they have forgotten the old ways.
  \item \textbf{Wolf-Pacted:} The clan has a standing truce with a local wolf pack. Wolves do not attack their herds. In exchange, the clan leaves every ninth calf unbranded — for the wolves. The pack sometimes fights beside them in battle (Cap 2, appears once per arc).
  \item \textbf{Salt-Misers:} This clan controls a small, bitter salt pan. They trade salt for iron with the Kuvani. They are wealthy but paranoid; any visitor is assumed to be a spy. +1 die to detect deception, but starting Position in parley is Desperate.
  \item \textbf{Curse-Bearers:} A broken oath from three generations ago still hangs over the clan. Their horses are prone to sickness, and their children are born with grey eyes. They seek a way to lift the curse — and will do almost anything to find it.
  \item \textbf{Road-Breakers:} This clan claims the right to open or close the high passes east. They carry a brass key (actually a bronze paiza) that Kuvani caravans recognise. They are isolationist, proud, and quick to draw steel over matters of passage.
\end{enumerate}

\paragraph*{Eastern Steppe Hazards (d6)}
When travelling through Eastern Ykrul territory, roll for a hazard (or replace a standard Clubs draw).
\begin{enumerate}
  \item \textbf{Frost-Rime Grass:} The grass is coated in a crystalline frost that cuts horse legs. Travel speed halved for one leg. On a failed Survival roll, one mount goes lame (mark 1 Fatigue or lose 1 Supply).
  \item \textbf{Wolf Howl at Dusk:} A pack shadows the party for a day and a night. They will not attack unless the party is weak or separated. If a character sleeps alone, roll Spirit+Resolve (DV 3) or be found by wolves at midnight.
  \item \textbf{Ice Lens Mirage:} The flat steppe plays tricks. Navigation rolls suffer –1 die for the leg. On a miss, the party wanders into a kurgan field — the dead are restless. Start a \textbf{Spirit Attention} [4] clock.
  \item \textbf{Salt Pan Crack:} A dried salt pan has a hidden fissure. A horse or wagon breaks through. Lose 1 Supply (salt pork ruined) and mark 1 Fatigue for the party as they pull the animal free.
  \item \textbf{Kuvani Border Markers:} The party has crossed into a Kuvani grazing ground without permission. Kuvani riders appear (Cap 2, non‑hostile unless provoked). They demand a story or a small offering of salt. Refusal starts a \textbf{Border Dispute} [4] clock.
  \item \textbf{Silent Sky:} No birds, no wind, no insect sounds. The Hollow is near. For the next scene, any roll that shows a 1 also summons a \textbf{Oath-Keeper} — it will appear at the end of the scene to ask one question. Answer poorly, and it takes a memory.
\end{enumerate}

\paragraph*{Eastern Ykrul Encounter (d8)}
When meeting an Eastern Ykrul camp or patrol, roll for their initial disposition and a twist.
\begin{enumerate}
  \item \textbf{Hungry Hosts:} The camp is low on supplies. They will trade — at ruinous prices. But they will also remember a fair deal. If you overpay (give 1 extra Supply), gain a \textbf{String} (Eastern Ykrul favour).
  \item \textbf{Grieving Clan:} A child died last night. The camp is in mourning. They will not trade or talk. If you offer a story about a lost loved one, they will share a meal. The story must be true.
  \item \textbf{Feud-Watch:} The clan is on alert for raiders from a rival banner. They mistake the party for scouts. The first social roll is at Desperate Position. A successful Deception or Bribe (cost 1 Boon or 1 Supply) clears the misunderstanding.
  \item \textbf{Salt-Broker’s Caravan:} A small caravan of salt traders, heavily armed. They are willing to trade salt for news of the west. They ask pointed questions about the Violet Steppe — they are planning a raid.
  \item \textbf{Kurgan Pilgrims:} A group of bone‑singers and their guards are travelling to a sacred kurgan. They will allow the party to accompany them if the party helps perform a ritual: each member must whisper a forgotten name into a horse skull. Afterwards, the party gains +1 die to resist fear for one day, but the Hollow knows one of your names.
  \item \textbf{Storm-Singer’s Challenge:} A sky‑speaker challenges the party to a weather‑calling contest. Roll Spirit+Lore vs. DV 4. On a success, the storm-singer is impressed and grants a \textbf{Sky’s Exception} (ignore one weather hazard). On a failure, the party is drenched and loses 1 Fatigue.
  \item \textbf{Hollow-Touched Camp:} The camp is silent. Everyone is sitting still, staring at a single point. A Oath-Keeper has passed through. If the party can speak the name of the broken oath, the camp will wake. Otherwise, they must leave an offering of salt and retreat.
  \item \textbf{Iron Grass Market:} A rare gathering of three clans at a neutral ford. The market lasts only one day. Goods for sale: remounts (cost 1 Supply each), salt (1 Supply per 2 portions), iron arrowheads (1 Boon for 1 die bonus on next ranged attack), and one cursed paiza (half price, but the previous owner’s ghost follows).
\end{enumerate}

\paragraph*{Eastern Ykrul Customs (d6)}
When interacting with an Eastern Ykrul camp, roll for a local custom that may affect social rolls.
\begin{enumerate}
  \item \textbf{No Left Hand:} Offering anything with the left hand is an insult. All social rolls using Presence+Sway suffer –1 die unless you explicitly use your right hand.
  \item \textbf{Fire First, Talk Later:} Guests must warm their hands at the central fire before speaking. If you refuse, the host will be offended — the first request you make is automatically refused.
  \item \textbf{The Empty Saddle:} One saddle in the camp is always empty. It belongs to a rider who died last winter. If you sit in it, you must tell a story about a death you witnessed. If the story is true, the ghost may give you a warning (GM reveals one future hazard). If false, the ghost haunts you for a day (you gain the \textbf{Haunted} condition).
  \item \textbf{No Naming the Dead:} Eastern Ykrul do not speak the names of the deceased for one year after death. If you accidentally name a dead person, the clan demands a blood-price: 1 Supply or a lock of your hair.
  \item \textbf{Salt Before Meat:} A meal is not blessed until the host throws a pinch of salt into the fire. If you eat before the salt, you are considered an oath‑breaker. The meal will not nourish you (you do not recover Fatigue from that rest).
  \item \textbf{The Ninth Breath:} When drinking fermented mare’s milk, you must pause after the eighth sip. If you take a ninth, you must immediately name an enemy. The clan will consider that enemy as their own for one season. The enemy may not appreciate this.
\end{enumerate}

\paragraph*{Eastern Ykrul Superstitions (burn for +1 die once)}
Add these to the standard Ykrul superstitions when in the eastern steppe.
\begin{itemize}
  \item \textbf{Iron to the Frost:} Press a cold iron nail into the ground before a white squall. Ignore the first \emph{Frost-Rime Grass} penalty.
  \item \textbf{Braid the Tail:} Tie a single knot in a horse’s tail hair. Cancel \emph{Wolf Howl at Dusk} or take –1 die on your next survival roll (your choice).
  \item \textbf{Name to the Hollow:} Whisper a dead enemy’s name into a salt pan at midnight. Gain +1 Effect when evading Oath-Keepers this scene, but mark \emph{Obligation} to the Eastern Khatun.
\end{itemize}

\paragraph*{Eastern Ykrul Boss Seed (d6)}
When the party faces a major antagonist from the Eastern Ykrul, roll for a unique foe.
\begin{enumerate}
  \item \textbf{The Frost-Khagan (Undying):} A chieftain who froze to death but did not stop riding. \emph{Phases}: Blizzard’s Call $\rightarrow$ Frozen Host $\rightarrow$ The Long Night. \emph{Cracks}: burn a horse tooth in a fire, speak his true name (carved on his saddle), or offer a living rider to take his place.
  \item \textbf{The Salt-Throat (Cursed Broker):} A salt merchant who sold his clan’s future for wealth. \emph{Phases}: Weigh the Debt $\rightarrow$ Salt Tithe $\rightarrow$ The Scales Close. \emph{Cracks}: pay a wergild in untainted salt, expose his false paiza, or cut his hostage string with a silver knife.
  \item \textbf{The Hollow’s Rider (Oath-Keeper in Flesh):} A warrior who made a pact with the Hollow to survive a winter famine. \emph{Phases}: Count the Breath $\rightarrow$ Ride the White $\rightarrow$ Claim the Ninth. \emph{Cracks}: feed him a memory of warmth, show him a red thread tied in eight knots, or outride him to the horizon — the Hollow cannot cross moving water.
  \item \textbf{The Iron Grandmother (Bone-Singer Revenant):} An ancient bone-singer who refused to die. \emph{Phases}: Song of Names $\rightarrow$ Kurgan Rising $\rightarrow$ The Dead Vote. \emph{Cracks}: sing a true name she does not know, return a stolen kurgan stone, or answer her riddle: “What weight is never paid?”
  \item \textbf{The Wolf-Banner (Pack-Leader):} A clan chief who has fused with a wolf spirit. \emph{Phases}: Stalk $\rightarrow$ Circle $\rightarrow$ Feast. \emph{Cracks}: break his banner, offer a raw heart to his pack, or survive a night without fire.
  \item \textbf{The Sky-Weeping Child (Curse Bearer):} A child born during a white squall, marked by the Hollow. \emph{Phases}: Winter Dreams $\rightarrow$ Frost-Fever $\rightarrow$ The Ninth Day. \emph{Cracks}: warm her hands with your own breath, speak her true name (her mother is still alive — she will tell you), or find the horse that was born at the same moment — it carries the other half of the curse.
\end{enumerate}

\paragraph*{Eastern Steppe Names (d8 each)}
For quick NPC generation, roll or choose:

\textbf{Male first names (d8):}  
1. Bortei, 2. Temurjin, 3. Khorchi, 4. Edil, 5. Sogata, 6. Toghrul, 7. Ogedei, 8. Chagatai.

\textbf{Female first names (d8):}  
1. Orghana, 2. Borte, 3. Altan, 4. Saruul, 5. Checheg, 6. Nergui, 7. Tolui, 8. Hulan.

\textbf{Clan epithets (d8):}  
1. Iron-Heart, 2. Frost-Mane, 3. Salt-Tongue, 4. Wolf-Brood, 5. Hollow-Sight, 6. Curse-Bearer, 7. Red-Hoof, 8. Silent-Banner.

\begin{tcolorbox}[colback=black!3,colframe=black!40!white,title={Eastern Ykrul at a Glance}]
\textbf{Core Traits:} Isolationist, pragmatic, harsher honour code, live among Kuvani but do not intermarry.  
\textbf{Key Resources:} Salt, iron, remounts, wolf truces, kurgan spirits.  
\textbf{Common Conflicts:} Feuds over salt pans, disputes with Kuvani over grazing, Hollow incursions in winter, curse‑lifting quests.  
\textbf{Tone:} Bleaker than western Ykrul. Fewer stories, more teeth. Hospitality is still sacred — but the guest is expected to earn it.  
\textbf{Ritual Objects:} Frost-bitten stirrups, bronze paizas, horse‑tooth chits, salt‑cakes, wolf skull standards.
\end{tcolorbox}